%--------------------------------------------------------------------------------
%
%
%--------------------------------------------------------------------------------
\pagebreak
\section*{AND Boolean Operation}
\addcontentsline{toc}{section}{AND Boolean Operation}

\instructionentry{Syntax}{
  \texttt{AND[.C/N] RegR, RegB, RegA} \\
  \texttt{AND[.C/N] RegR, RegB, SignedImmed15}
  
  \texttt{AND[.C/N/D/W/H/B] RegR, SignedImmed13( RegB )} \\
  \texttt{AND[.C/N/D/W/H/B] RegR, RegX ( RegB )}
}

\instructionentry{Format}{

	The AND instruction uses the register, the immediate-19, the indexed and the register indexed instruction formats. \\
	
	\begin{flushleft}	

	\drawInstrFormat
  	{0,1,3,5,5.5,6,6.5,8.5,9.5,11.5,16}
  	{\OpCodeGrpALU,\OpCodeAND, RegR, N, C, 0, RegB,  RegA, 0}
	
	\drawInstrFormat
  	{0,1,3,5,5.5,6,6.5,8.5,16}
  	{\OpCodeGrpALU,\OpCodeAND, RegR, N, C, 1, RegB, val}

	\drawInstrFormat
  	{0,1,3,5,5.5,6,6.5,8.5,9.5,16}
  	{\OpCodeGrpMEM,\OpCodeAND, RegR, N, C, 0, RegB, dw, ofs}
 
	\drawInstrFormat
  	{0,1,3,5,5.5,6,6.5,8.5,9.5,11.5,16}
  	{\OpCodeGrpMEM,\OpCodeAND, RegR, N, C, 1, RegB, dw, RegX, 0}

	\end{flushleft}
}

\instructionentry{Description}{

The \texttt{AND} instruction performs the binary AND operation on two 64-bit operands and stores the result to the target register \texttt{RegR}. Upon input the general register \texttt{regB} can be negated. The result can also be negated by setting the "C" bit. The result can be negated by setting the "N" bit. The instruction supports the immediate, the register and the two memory reference instruction formats. The boolean operation itself does not cause any traps.The indexed formats may cause a memory reference associated trap.
}

\instructionentry{Operation}{

	\texttt{if ( instr.[C] ) RegB <- not( RegB );} \\
	\texttt{RegR <- RegB \& RegA;}           (register format)\\
	\texttt{RegR <- RegB \& immOperand();}   (immediate format)\\
	\texttt{RegR <- RegR \& memOperand();}   (indexed formats)\\
	\texttt{if ( instr.[N] ) RegR <- not( RegR );}
}

\instructionentry{Exceptions}{

	Data Alignment Trap\\
	Memory Access Violation Trap\\
	Memory Protection Violation Trap\\
	Data TLB Miss Trap
}

\instructionentry{Notes}{

	None. 
}




